% TD 3 : Traefik et Portainer
\td{3}{Traefik et Portainer}

Dans ce TD, j'ai placé Traefik, un reverse proxy, devant mon cluster Swarm : chaque application répond à son propre nom de domaine au lieu d'un port.
J'ai ensuite installé Portainer pour gérer le cluster depuis le navigateur.

% =====================================================================
\exercice{1}{Déployer Traefik}

\subsection{Relancer le cluster}

J'ai repris le cluster du TD~2 en publiant le port 80 du manager, pour que le navigateur du PC puisse joindre Traefik.

\begin{fichier}{mfp/swarm/compose.yml}
  manager:
    image: docker:dind
    hostname: manager
    privileged: true
    ports:
      - "80:80" # Traefik, reachable from the host browser
\end{fichier}

\begin{commandes}[Depuis mfp/swarm]
docker compose up -d
docker compose exec manager docker node ls
\end{commandes}

Seul le manager a été recréé, et \code{node ls} affichait de nouveau les 4 nœuds \code{Ready}.

\subsection{Déployer Traefik sur le manager}

La stack contient Traefik, placé sur le manager, et \code{whoami}, une application de test qui se déclare à Traefik par ses labels.

\begin{fichier}{mfp/swarm/traefik.compose.yml}
  # Deploy the Whoami application
  whoami:
    image: traefik/whoami
    networks:
      - web
    deploy:
      labels:
        - "traefik.enable=true"
        - "traefik.http.routers.whoami.rule=Host(`whoami.swarm.localhost`)"
        - "traefik.http.routers.whoami.entrypoints=web"
        - "traefik.http.services.whoami.loadbalancer.server.port=80"

# Define the overlay network for Swarm
networks:
  web:
    external: true
\end{fichier}

\needspace{4cm}
Le réseau \code{web} est déclaré \code{external} : je l'ai donc créé avant le déploiement, en mode overlay pour qu'il couvre tous les nœuds.

\begin{commandes}
docker compose exec manager docker network create --driver overlay web
docker compose exec manager docker network ls --filter name=web
\end{commandes}

J'ai ensuite supprimé la stack et le service du TD~2, puis j'ai déployé la stack Traefik.

\begin{commandes}
docker compose exec manager docker stack rm hello
docker compose exec manager docker service rm web
docker compose cp traefik.compose.yml manager:/home/manager/
docker compose exec manager docker stack deploy -c /home/manager/traefik.compose.yml traefik
docker compose exec manager docker service ls
\end{commandes}

\capture[0.22\textheight]{td3-03-traefik-deploy.png}{Les services \code{traefik_traefik} et \code{traefik_whoami} sont à 1/1. Seul Traefik publie un port (80) ; whoami reste derrière lui.}

\subsection{Accéder à Traefik depuis le navigateur}

Aucun vrai domaine ne pointe vers mon PC : j'ai donc ajouté les 2 adresses dans le fichier hosts de Windows, ouvert en administrateur.

\begin{fichier}[text]{C:\textbackslash Windows\textbackslash System32\textbackslash drivers\textbackslash etc\textbackslash hosts}
127.0.0.1   traefik.swarm.localhost
127.0.0.1   whoami.swarm.localhost
\end{fichier}

\capture[0.27\textheight]{td3-05-traefik-dashboard.png}{Le tableau de bord sur \code{traefik.swarm.localhost}. Le point d'entrée \code{web} écoute sur le port 80 et les 4 routeurs sont sans erreur.}

\capture[0.34\textheight]{td3-06-whoami.png}{La page \code{whoami.swarm.localhost}. L'en-tête \code{Host} et les en-têtes \code{X-Forwarded-*} ajoutés par Traefik montrent que la requête est passée par le reverse proxy.}

% =====================================================================
\exercice{2}{Déployer une autre application}

\subsection{Tester l'application avec Docker Compose}

L'application de vote compte 5 services : \code{vote}, \code{redis}, \code{worker}, \code{db} et \code{result}.
Je l'ai d'abord lancée en local, avec le vote sur le port 8080 et les résultats sur le port 8081.

\begin{commandes}[Depuis voting-app]
docker compose up -d
docker compose ps --format 'table {{.Name}}\t{{.Image}}\t{{.Status}}\t{{.Ports}}'
\end{commandes}

\probleme{La première tentative s'est arrêtée sur \code{429 Too Many Requests}, car Docker Hub limite les téléchargements sans compte.}{Je me suis connecté avec \code{docker login}, puis j'ai relancé \code{docker compose up -d}.}

J'ai envoyé 5 votes avec \code{curl}, ce qui revient à cliquer sur les boutons du site.

\begin{commandes}
curl -X POST -d vote=a http://localhost:8080/   # Cats, 3 fois
curl -X POST -d vote=b http://localhost:8080/   # Dogs, 2 fois
\end{commandes}

\capture[0.21\textheight]{td3-09-result-local.png}{Le site des résultats sur \code{localhost:8081} : 5 votes, 60~\% pour Cats et 40~\% pour Dogs.}

Une fois le test terminé, j'ai arrêté l'application avec \code{docker compose down}.

\subsection{Déployer l'application sur le cluster Swarm}

La stack \code{docker-stack.yml} utilise des images déjà publiées, car Swarm ne construit pas d'images, et demande 2 répliques de \code{vote} et de \code{worker}.
Je me suis connecté à Docker Hub dans le manager, et \code{--with-registry-auth} transmet cette connexion aux autres nœuds.

\begin{commandes}[Depuis mfp/swarm]
docker compose exec manager docker login
docker compose cp ../../voting-app/docker-stack.yml manager:/home/manager/
docker compose exec manager docker stack deploy --with-registry-auth -c /home/manager/docker-stack.yml voting
docker compose exec manager docker stack services voting
\end{commandes}

Au bout d'une minute environ, les 5 services étaient au complet, mais les ports 8080 et 8081 n'étaient publiés qu'à l'intérieur du cluster.

\subsection{Publier le vote et les résultats avec Traefik}

\subsubsection{Labels et redéploiement}

Dans une nouvelle version de la stack, j'ai branché \code{vote} et \code{result} sur le réseau \code{web} et remplacé leurs ports par des labels Traefik.
Voici la partie de \code{vote} ; celle de \code{result} est identique, avec \code{result.swarm.localhost}.

\begin{fichier}{voting-app/docker-stack.traefik.yml}
    networks:
      - frontend
      - web
    deploy:
      replicas: 2
      labels:
        - "traefik.enable=true"
        - "traefik.http.routers.vote.rule=Host(`vote.swarm.localhost`)"
        - "traefik.http.routers.vote.entrypoints=web"
        - "traefik.http.services.vote.loadbalancer.server.port=80"
\end{fichier}

\needspace{4cm}
J'ai redéployé sous le même nom, \code{voting}, pour que Swarm mette à jour les 5 services existants.

\begin{commandes}
docker compose cp ../../voting-app/docker-stack.traefik.yml manager:/home/manager/
docker compose exec manager docker stack deploy --with-registry-auth -c /home/manager/docker-stack.traefik.yml voting
docker compose exec manager docker stack services voting
\end{commandes}

\subsubsection{Vérification}

J'ai ajouté les 2 nouvelles adresses dans le fichier hosts.

\begin{fichier}[text]{hosts}
127.0.0.1   vote.swarm.localhost
127.0.0.1   result.swarm.localhost
\end{fichier}

\capture[0.30\textheight]{td3-15-traefik-routers.png}{La page HTTP Routers de Traefik. Les routeurs \code{vote@swarm} et \code{result@swarm} ont été trouvés grâce aux labels et sont en succès.}

J'ai de nouveau envoyé 5 votes, cette fois 2 pour Cats et 3 pour Dogs.

\begin{commandes}
curl -X POST -d vote=a http://vote.swarm.localhost/   # Cats, 2 fois
curl -X POST -d vote=b http://vote.swarm.localhost/   # Dogs, 3 fois
\end{commandes}

\capture[0.21\textheight]{td3-14-result-traefik.png}{Le site des résultats sur \code{result.swarm.localhost} : 40~\% pour Cats et 60~\% pour Dogs. La base du cluster ne contient que les votes envoyés au cluster.}

% =====================================================================
\exercice{3}{Déployer Portainer}

\subsection{Déployer Portainer sur le cluster}

J'ai téléchargé la stack de Portainer pour Swarm : un agent sur chaque nœud (mode \code{global}) et l'interface sur le manager.

\begin{commandes}[Depuis mfp/swarm]
curl -sL https://downloads.portainer.io/ce-lts/portainer-agent-stack.yml -o portainer-agent-stack.yml
docker compose cp portainer-agent-stack.yml manager:/home/manager/
docker compose exec manager docker stack deploy --with-registry-auth -c /home/manager/portainer-agent-stack.yml portainer
docker compose exec manager docker stack services portainer
docker compose exec manager docker ps --filter name=portainer --format 'table {{.Names}}\t{{.Image}}\t{{.Status}}'
\end{commandes}

\capture[0.30\textheight]{td3-16-portainer-deploy.png}{\code{portainer_agent} est à 4/4, un agent par nœud, et \code{portainer_portainer} est à 1/1. Les 2 conteneurs tournent bien sur le manager.}

\subsection{Publier Portainer avec Traefik}

Comme pour \code{whoami}, j'ai branché Portainer sur le réseau \code{web} et ajouté ses labels, vers son port 9000.

\begin{fichier}{mfp/swarm/portainer-agent-stack.traefik.yml}
    networks:
      - agent_network
      - web
    deploy:
      mode: replicated
      replicas: 1
      placement:
        constraints: [node.role == manager]
      labels:
        - "traefik.enable=true"
        - "traefik.http.routers.portainer.rule=Host(`portainer.swarm.localhost`)"
        - "traefik.http.routers.portainer.entrypoints=web"
        - "traefik.http.services.portainer.loadbalancer.server.port=9000"

networks:
  agent_network:
    driver: overlay
    attachable: true
  web:
    external: true
\end{fichier}

J'ai redéployé sous le même nom de stack, ce qui a mis à jour les 2 services sans les recréer.

\begin{commandes}
docker compose cp portainer-agent-stack.traefik.yml manager:/home/manager/
docker compose exec manager docker stack deploy --with-registry-auth -c /home/manager/portainer-agent-stack.traefik.yml portainer
docker compose exec manager docker stack services portainer
\end{commandes}

\subsection{Ouvrir Portainer dans le navigateur}

\subsubsection{Setup token}

J'ai ajouté l'adresse de Portainer dans le fichier hosts.

\begin{fichier}[text]{hosts}
127.0.0.1   portainer.swarm.localhost
\end{fichier}

Au premier lancement, Portainer demande un setup token, que j'ai trouvé dans ses logs.

\begin{commandes}
docker compose exec manager docker service logs --raw portainer_portainer 2>&1 | grep -A2 setup_token=
\end{commandes}

\subsubsection{Création du compte}

Sur \url{http://portainer.swarm.localhost}, j'ai créé le compte \code{admin} avec ce token et un mot de passe d'au moins 12 caractères.
Portainer ne laisse que 5 minutes après son démarrage pour le faire.

\capture{td3-21-portainer-home.png}{L'accueil de Portainer. L'environnement \code{primary} est \code{Up}, de type Swarm, avec les 4 nœuds trouvés par les agents.}

% =====================================================================
\exercice{4}{Redéployer depuis Portainer}

\subsection{Supprimer la stack depuis le manager}

J'ai supprimé la stack de l'application de vote, créée en ligne de commande.

\begin{commandes}[Depuis mfp/swarm]
docker compose exec manager docker stack ls
docker compose exec manager docker stack rm voting
docker compose exec manager docker stack ls
\end{commandes}

\code{stack rm} a supprimé les 5 services et les 2 réseaux de \code{voting} : il ne restait plus que \code{portainer} et \code{traefik}.

\subsection{Republier la stack depuis Portainer}

Dans Portainer, j'ai ouvert \code{primary}, puis \og Stacks\fg{} et \og Add stack\fg{}.
J'ai nommé la stack \code{voting}, envoyé le fichier \code{docker-stack.traefik.yml} de la question 2.3 avec la méthode \og Upload\fg{}, puis cliqué sur \og Deploy the stack\fg{}.

\capture[0.30\textheight]{td3-24-portainer-add-stack.png}{Le formulaire de création de la stack \code{voting}. Portainer fait l'équivalent d'un \code{docker stack deploy}.}

\capture[0.30\textheight]{td3-25-portainer-stacks.png}{\code{voting} est de retour, avec un contrôle \og Total\fg{} et créée par \code{admin}. Les stacks créées en ligne de commande restent \og Limited\fg{} : Portainer ne peut pas modifier leur fichier.}

\subsection{Vérifier l'accès depuis le navigateur}

Après ce redéploiement, le vote et les résultats répondaient de nouveau à leurs adresses : la stack déployée par Portainer fonctionne comme celle déployée en ligne de commande.

\probleme{La page des résultats affichait \og No votes yet\fg{} : la base avait été recréée sur \code{node2} au lieu de \code{node3}, avec un volume neuf.}{Je n'ai rien corrigé : les anciens votes sont restés dans le volume de \code{node3}.}

\subsection{Déployer l'API de MFP depuis Portainer}

\subsubsection{Création de la stack}

Ma stack reprend la base et le serveur de MFP. Le serveur utilise l'image publiée sur GHCR et passe par Traefik sur \code{api.swarm.localhost}.

\begin{fichier}{mfp/swarm/mfp-api.stack.yml}
  server:
    image: ghcr.io/dark-shad0w/clusterisation-conteneurs/server:latest
    networks:
      - backend
      - web
    deploy:
      labels:
        - "traefik.enable=true"
        - "traefik.http.routers.mfp-api.rule=Host(`api.swarm.localhost`)"
        - "traefik.http.routers.mfp-api.entrypoints=web"
        - "traefik.http.services.mfp-api.loadbalancer.server.port=3000"

networks:
  backend:
  # overlay network shared with Traefik (created in TD3 exercise 1)
  web:
    external: true
\end{fichier}

Dans Portainer, j'ai créé la stack \code{mfp} comme à la question 4.2, avec le fichier \code{mfp-api.stack.yml}.

\capture[0.22\textheight]{td3-31-portainer-mfp-services.png}{Les 2 services de la stack \code{mfp} sont au complet. \code{mfp_server} utilise l'image publiée sur GHCR.}

\probleme{Au premier démarrage, le serveur s'arrêtait sur \code{getaddrinfo ENOTFOUND db}, car Swarm lance la base et le serveur en même temps.}{Je n'ai rien eu à faire : Swarm a relancé le serveur, qui s'est connecté à la base au deuxième essai.}

\subsubsection{Test avec Bruno}

J'ai ajouté l'adresse de l'API dans le fichier hosts.

\begin{fichier}[text]{hosts}
127.0.0.1   api.swarm.localhost
\end{fichier}

\needspace{4cm}
Pour ne pas modifier ma collection Bruno, j'en ai fait une copie qui vise \code{http://api.swarm.localhost} au lieu de \code{http://localhost:3000}.

\begin{commandes}[Depuis mfp]
cp -r mfp-bruno-api-test-collection ~/mfp-bruno-swarm
sed -i 's#http://localhost:3000#http://api.swarm.localhost#' ~/mfp-bruno-swarm/*.yml
cd ~/mfp-bruno-swarm && bru run
\end{commandes}

\capture[0.34\textheight]{td3-33-bruno-swarm.png}{Les 3 requêtes répondent 200 OK et les 3 tests passent. L'API tourne sur le cluster et répond depuis le PC sur \code{api.swarm.localhost}.}
